\documentclass[11pt]{article}
\usepackage[margin=1in]{geometry}
\usepackage{amsmath,amssymb,amsthm,mathtools}
\usepackage[T1]{fontenc}
\usepackage[colorlinks=true,linkcolor=blue,citecolor=blue,urlcolor=blue]{hyperref}
\usepackage{microtype}
\newtheorem{theorem}{Theorem}
\newtheorem{lemma}[theorem]{Lemma}
\newtheorem{corollary}[theorem]{Corollary}
\newtheorem{proposition}[theorem]{Proposition}
\theoremstyle{remark}
\newtheorem{remark}[theorem]{Remark}
\newcommand{\Q}{\mathbb Q}
\newcommand{\E}{\mathbb E}
\newcommand{\TV}{d_{\mathrm{TV}}}
\newcommand{\tr}{\operatorname{tr}}
\newcommand{\supp}{\operatorname{supp}}
\newcommand{\Pf}{\operatorname{Pf}}
\newcommand{\ket}[1]{\lvert #1\rangle}
\newcommand{\bra}[1]{\langle #1\rvert}
\title{Coherent preparation and observable access for\newline partition-projected Slater determinants}
\author{Conditional transfer note; language-model-assisted research draft}
\date{7 October 2026}
\begin{document}
\maketitle
\begin{abstract}
Assuming the homogeneous coefficient FPRAS of Chen and Liu, we give a uniform finite-bit compiler for an explicitly supplied rational-complex Slater determinant projected onto arbitrary disjoint occupation quotas. The compiler prepares the normalized projected state in polynomial time, with no inverse projection-success factor. Its randomness is fixed per depth and shared across all prefixes. A probability-weighted unitary hybrid bound avoids an exponential union bound and requires relative counting precision proportional to the target trace error, retaining an inverse-square counting-accuracy dependence. Exact support tests, coefficient bounds, conditional oracles, phase computation, and reversible workspace are acquired from the input. Classical sparse-observable and fixed-copy purity estimates follow with a clipping argument. Counting and conditional state preparation are established primitives: the contribution is their explicit composition for this physical input class. The imported FPRAS has not been independently certified in full or executed here; global novelty remains unverified.
\end{abstract}

\section{Input, conclusion, and imported result}
Fix an increasing order of $m$ fermion modes. The input consists of a full-column-rank matrix $A\in\Q(i)^{m\times r}$ and a partition $[m]=E_1\sqcup\cdots\sqcup E_k$, with integers $0\le q_j\le |E_j|$ and $\sum_jq_j=r$. All rational numerators and denominators have finite binary encodings. Write $L$ for their total input length. For an $r$-element subset $S$, let $A_S$ have rows in increasing order and define
\[
 a(S)=\det A_S,\qquad
 \mathcal F=\{S:|S\cap E_j|=q_j\ \forall j\},\qquad
 Z=\sum_{S\in\mathcal F}|a(S)|^2.
\]
Amplitudes are zero on other occupation strings. When $Z>0$, the target state and Born law are
\begin{equation}\label{eq:target}
 \ket\psi=Z^{-1/2}\sum_{S\in\mathcal F}a(S)\ket S,
 \qquad p(S)=|a(S)|^2/Z.
\end{equation}
Trace distance means $\frac12\|\rho-\sigma\|_1$. Input acquisition and optimization of $A$ are outside this interface.

The only fresh load-bearing import is Chen--Liu \cite{CL}, Theorem~5.1 (printed pp.~19--20), with the rational implementation in Remark~A.4 (p.~38). In the form used here it states: for nonnegative homogeneous multi-affine log-concave weight polynomials $f,g$ of degrees $a,b$ with $a+b=m$, coefficient-value oracles, positive initializing supports, and supplied positive-coefficient range bounds $R_f,R_g$, there is an FPRAS for
\begin{equation}\label{eq:import}
 [x_1\cdots x_m]fg
\end{equation}
with success probability at least $3/4$ and, on \emph{every execution},
\begin{equation}\label{eq:cost}
 O\!\left(m^{10}B^2\eta^{-2}\log^2(2B/\eta)\right),
 \qquad B=m+\log(R_fR_g),
\end{equation}
arithmetic operations and oracle calls for relative error $\eta$. Polynomial-time rational oracles returning polynomial-length rationals give polynomial bit time on every execution. The source implements rational choices with capped fair-coin segments; aborts are counted in its failure probability. Median amplification gives failure $\beta$ with an additional $O(\log(1/\beta))$ factor. Its zero target always produces zero. We nevertheless acquire an exact support test independently.

\begin{theorem}[Conditional projected-Slater access]\label{thm:main}
Assume the imported result above. For the explicit input just specified, one can decide $Z=0$ exactly in polynomial bit time. If $Z>0$, there are uniform randomized polynomial-bit algorithms for relative approximation of $Z$, sampling $p$ to arbitrary total-variation error $\tau>0$, and the following coherent task. Given rational $0<\epsilon,\delta<1$, a randomized classical compiler outputs a polynomial-size quantum circuit whose output has trace distance at most $\epsilon$ from $\ket\psi$, with probability at least $1-\delta$ over compiler randomness. Compilation and circuit size are polynomial in $L,m,\epsilon^{-1},\log(\delta^{-1})$. No inverse power of $Z$ or of the normalized projection probability occurs.
\end{theorem}

For $m=2n,r=n$, $E_j=\{j\uparrow,j\downarrow\}$ and $q_j=1$, this is single-occupancy Gutzwiller projection of any number-conserving, possibly spin-mixed, Slater state. Two independent spin Slaters are a special case. Any number of flavors with exactly one mode per site is also covered. The normalized unprojected Slater has squared norm $\det(A^\dagger A)$; hence its projection probability is $Z/\det(A^\dagger A)$ and can be relatively approximated as well. Full rank does not imply $Z>0$: on two sites, occupied orbitals localized to the two spin modes of the first site give a full-rank input whose single-occupancy projection vanishes.

\section{Acquiring the coefficient and prefix interfaces}
\begin{lemma}[Determinant coefficient identity]\label{lem:det}
The polynomials
\begin{equation}\label{eq:fg}
 f_A(x)=\det(A^\dagger\operatorname{diag}(x)A),\qquad
 g(x)=\prod_{j=1}^k e_{|E_j|-q_j}(x_{E_j})
\end{equation}
satisfy the hypotheses of \eqref{eq:import}, and $[x_1\cdots x_m]f_Ag=Z$.
\end{lemma}
\begin{proof}
Cauchy--Binet gives
\[
 f_A(x)=\sum_{|S|=r}|\det A_S|^2x^S.
\]
For real $x>0$, $A^\dagger\operatorname{diag}(x)A$ is positive definite Hermitian and affine in $x$. Concavity of $\log\det$ proves log-concavity of $f_A$. It is multi-affine and homogeneous of degree $r$, with nonnegative rational coefficients. The elementary symmetric polynomial $e_t$ is homogeneous real stable and therefore log-concave; equivalently one may use the standard matroid basis-polynomial result \cite{BH,ALOV}. Products in disjoint variable blocks preserve log-concavity, so $g$ has degree $m-r$. Its coefficient at $x^T$ is one precisely when $[m]\setminus T\in\mathcal F$. The full coefficient is thus the stated sum. In the single-occupancy spin case, $g=\prod_j(x_{j\uparrow}+x_{j\downarrow})$, making its log-concavity immediate.
\end{proof}

Both coefficient oracles are acquired: one computes an exact squared rational-complex minor or a quota membership test. Gaussian elimination supplies an initial nonzero minor, and blockwise choices supply an initial support of $g$. To acquire coefficient bounds, choose a common denominator $D$ clearing the real and imaginary parts of all entries. Then $DA$ is Gaussian integral. If $\det A_S\ne0$, its squared modulus is at least $D^{-2r}$. A permutation or Hadamard bound supplies a maximum with polynomial logarithm in $L,m$. Consequently $\log R_f$ is polynomial and $R_g=1$. Exact determinant computation has polynomial bit cost; no minimum singular-value promise is supplied or needed.

For a prefix of occupation bits, let $I$ be forced-in modes and $J$ forced-out modes. Replace the coefficient of $f_A$ by
\[
 |\det A_S|^2\,\mathbf1\{I\subseteq S,\ S\cap J=\varnothing\};
\]
restrict $g$ complementarily, forcing $J$ in and $I$ out. These operations are coordinate differentiation, exclusion, and multiplication by forced variables. Keeping the original ground set and forced variables retains degrees $r,m-r$, exactly as in the source's conditional construction (p.~24). The derivative closure follows from multi-affinity and the limit
\[
 \partial_i f(x)=\lim_{t\to\infty}t^{-1}f(x+t e_i),
\]
followed by iteration; exclusion follows by a boundary limit, and forced-variable factors add concave logarithms. Nonzero restrictions remain log-concave and their coefficient ranges do not increase. Each restricted query is one original query plus polynomial membership work.

The support of nonzero minors is a represented linear matroid; the quota family is the bases of a partition matroid. Polynomial-time matroid intersection \cite{Edmonds}, with exact rational rank queries, decides whether a common base exists and returns one. Contracting $I$ and deleting $J$ gives the prefix test, with immediate rejection for incompatible quotas or dependent forced elements. A common witness supplies positive initializing supports for both restricted polynomials. Thus every child count, including exact detection of an empty child subtree, comes from a single uniform polynomial-time program on the prefix. There is no exponentially stored conditional decision tree.

This proves the counting assertion. For classical sampling, use fresh randomness at the visited prefixes, normalized nonnegative child estimates, and a known feasible fallback if all estimates vanish. Relative error $\eta$ on both children changes their conditional law by at most $2\eta$; either failure has probability at most $2\beta$. A sequential coupling therefore bounds the final total-variation error by $2m\eta+2m\beta$. Choosing both parameters at most $\tau/(8m)$ leaves room for capped rational sampling. This also proves that the classical sampler is acquired, rather than assumed.

\section{A weighted-prefix coherent error bound}
We state the compilation step separately because an exponential union bound over prefix counts would destroy it. Grover--Rudolph \cite{GR} already use recursive conditional probabilities for preparation and explicitly allow probabilistic integration. The following lemma supplies quantitative finite-error accounting for a fixed shared tape.

\begin{lemma}[Randomized conditional-count compiler]\label{lem:compiler}
Let $p$ be a distribution on $m$ bits. Suppose exact prefix feasibility and two-child relative counts have polynomial-time, finite-bit implementations with a fixed worst-case time bound. If each child count has relative error $\eta\le1/2$ except with probability $\beta$, a randomized choice of tapes gives controlled rotations whose positive-amplitude output $\ket{\phi_s}$ satisfies, with probability at least $1-\delta$,
\begin{equation}\label{eq:hybridbound}
 \left\|\ket{\phi_s}-\sum_x\sqrt{p(x)}\ket x\right\|_2
 \le 4m\eta+2m\sqrt{\beta/\delta}.
\end{equation}
The tapes are shared across all prefixes of a depth; no independence between their errors is required.
\end{lemma}
\begin{proof}
Choose one random tape $s_d$ for each depth $d=0,\ldots,m-1$, split into fixed padded independent segments for both child counts and each median repetition. For each fixed feasible prefix $h$, the usual FPRAS marginal guarantee still holds. Force an empty child's estimate to zero, normalize nonnegative estimates to $q_s(b\mid h)$, and use a deterministic feasible branch if normalization vanishes. At infeasible prefixes choose a legal sentinel rotation. These definitions are total, even on erroneous tapes.

Call $h$ bad when either child estimate fails and put
\[
 B_d(s)=\sum_{|h|=d}p(h)\mathbf1\{h\text{ is bad}\},\qquad
 B(s)=\sum_{d=0}^{m-1}B_d(s).
\]
The true prefix probabilities do not depend on the tapes. Linearity of expectation gives $\E B\le2m\beta$, without a union bound or prefix independence. At a good prefix, for a nonzero child probability,
\[
 \frac{1-\eta}{1+\eta}\le\frac{q_s(b\mid h)}{p(b\mid h)}
 \le\frac{1+\eta}{1-\eta},\qquad
 \sum_b\bigl(\sqrt{q_s(b\mid h)}-\sqrt{p(b\mid h)}\bigr)^2\le16\eta^2.
\]
Zero children are treated exactly. At a bad prefix the last squared distance is at most two.

Let $T_d$ be the true controlled rotation at stage $d$ and $G_d$ its fixed-tape counterpart, extended to unitaries by real two-dimensional rotations on the next bit. The identity
\[
 G_m\cdots G_1-T_m\cdots T_1
 =\sum_{d=1}^mG_m\cdots G_{d+1}(G_d-T_d)T_{d-1}\cdots T_1
\]
applied to $\ket{0^m}$ has the \emph{true} prefix amplitudes at the altered stage. Subsequent $G$'s preserve norm. The $d$th term is therefore at most $\sqrt{16\eta^2+2B_{d-1}}$. The triangle and Cauchy--Schwarz inequalities give
\[
 \|\phi_s-\phi_p\|_2\le4m\eta+\sqrt{2mB(s)}.
\]
Markov's inequality gives $B(s)\le2m\beta/\delta$ with probability at least $1-\delta$, proving \eqref{eq:hybridbound}.
\end{proof}

Choose
\begin{equation}\label{eq:parameters}
 \eta=\frac{\epsilon}{16m},\qquad
 \beta=\frac{\delta\epsilon^2}{64m^2}.
\end{equation}
The positive-amplitude vector error is at most $\epsilon/2$ with probability $1-\delta$. Amplification costs $\log(1/\beta)$, not $\beta^{-1}$. Combining $2m$ counts with \eqref{eq:cost} gives an arithmetic/oracle bound
\[
 O\!\left(m^{13}B^2\epsilon^{-2}
 \log^2(2Bm/\epsilon)\log\frac{m}{\delta\epsilon}\right),
\]
up to constants and the polynomial bit cost of acquired oracles. In particular, converting a total-variation estimate of order $\epsilon^2$ into fidelity is unnecessary: the direct hybrid retains an inverse-square accuracy dependence.

\section{Reversible implementation and determinant phases}
Fixing the tapes turns each count program into a deterministic finite-bit computation. Its source-guaranteed time cap applies on \emph{every} prefix and tape, which is essential to one uniform quantum circuit. Reversible simulation can retain computational history, copy the rational conditional estimate, and reverse the computation. The history is polynomial in the time bound; no claim of logarithmic workspace is needed. Compute $q_s(1\mid h)$ reversibly, apply the corresponding controlled rotation, and uncompute all count and arithmetic workspace before the next stage. Fixed tapes are parameters of the circuit description. There are $m$ copies of a uniform prefix program, rather than exponentially many hardwired counts.

Rational square roots and rotations can be approximated to uniform additive angle, hence operator, error $O(\epsilon/m)$. This includes $q=0,1$ and arbitrarily small positive branch probabilities. Standard reversible finite-bit arithmetic and gate synthesis yield polynomial circuits; the sum of rotation errors can be made at most $\epsilon/4$. This is an additive operator-accuracy requirement, not relative accuracy in a tiny amplitude.

At the end, compute $a(S)=\det A_S$ exactly. Define the phase function as $a(S)/|a(S)|$ on nonzero feasible minors and as one on every other string. The latter includes wrong cardinality and impossible prefixes: the implemented function must be total under synthesis leakage. Gaussian-rational determinant arithmetic has polynomial output length. For the phase, select the larger-magnitude real or imaginary component and compute the angle in a bounded-ratio quadrant; standard finite-bit elementary-function approximation then requires polynomial precision. The acquired nonzero-minor bit bound also excludes an uncharged zero-discrimination assumption. Apply the phase diagonally and uncompute its workspace, with operator error at most $\epsilon/4$.

Exact phases preserve the vector error of Lemma~\ref{lem:compiler}. The determinant in increasing mode order already includes the fermionic sign, so complex phases are immaterial to diagonal sampling but essential to coherent reconstruction and off-diagonal observables. For unit vectors, pure-state trace distance is no larger than their vector distance. Adding the $\epsilon/2$ counting error and $\epsilon/2$ synthesis budget proves Theorem~\ref{thm:main}. The success guarantee is over compiler seeds; the procedure does not certify which particular seed succeeds. Efficient reversible access to the coefficient oracle is indispensable: an opaque classical matroid oracle alone would not justify this coherent conclusion.

\section{Sparse observables with finite sampling error}
Let $O$ be a Hermitian operator with an explicit polynomial-time row enumerator, polynomially many nonzero entries per row, polynomial-length rational-complex entries, and supplied bound $\|O\|\le M$. On a sampled $S$ with $a(S)\ne0$, compute
\[
 X(S)=\operatorname{Re}\frac{(Oa)(S)}{a(S)}.
\]
The row enumerator and exact minor queries suffice; amplitudes on excluded strings are zero. Then
\begin{equation}\label{eq:localmoments}
 \E_pX=\bra\psi O\ket\psi,\qquad
 \E_pX^2\le\frac1Z\sum_{a(S)\ne0}|(Oa)(S)|^2\le M^2.
\end{equation}
The local estimator is established prior art; sample and amplitude-ratio access of this sort is studied by Havl\'{i}\v{c}ek \cite{Hav}. Large pointwise ratios still have polynomial rational encoding length, but total-variation sampling alone does not control an unbounded estimator.

\begin{corollary}[Acquired classical observable estimates]\label{cor:observable}
Under the import, sparse-observable expectations have additive error $u$ and failure $\delta$ in time polynomial in $L,m,M/u,\log(\delta^{-1})$ and the row-enumerator cost. There is no inverse $Z$ dependence.
\end{corollary}
\begin{proof}
If $u\ge M$, zero suffices; otherwise take
\[
 K=4M^2/u,\qquad \tau=u^2/(32M^2),\qquad
 X_K=\max(-K,\min(K,X)).
\]
Equation~\eqref{eq:localmoments} gives clipping bias at most $M^2/K=u/4$. Sampling from a law $q$ with $\TV(p,q)\le\tau$ adds bias at most $2K\tau=u/4$. Moreover
\[
 \E_qX_K^2\le M^2+2K^2\tau\le2M^2.
\]
Independent median-of-means with $O(M^2u^{-2}\log(\delta^{-1}))$ samples estimates the clipped mean within $u/2$. All sampler and arithmetic tolerances are charged, completing the claim.
\end{proof}

The same proof applies to any fixed number of copies by taking the product input law and amplitude ratios. A subsystem SWAP is norm one and one-sparse, so arbitrary-subsystem purity $\tr(\rho_B^2)$ has a polynomial additive classical estimate. Fixed-order cyclic permutations similarly estimate R\'{e}nyi moments. Turning an exponentially small moment into additive R\'{e}nyi entropy requires relative precision and can cost exponentially much. Sparse Hamiltonian energies and variances, correlations, and polynomially sized local reduced density matrices are accessible; useful ground-state overlap, gap certificates, and optimization are not acquired.

\section{Comparison, boundaries, and verification status}
Chen--Liu Corollary~C.13(a) already specializes the counting primitive to determinant weights restricted to matroid bases. Counting the norm here is therefore a direct application of published prior art. Recursive coherent conditional preparation is Grover--Rudolph's method, including their probabilistic-integration observation. Our explicit weighted-prefix bound and acquisition of phases and finite-bit oracles assemble these ingredients into the projected-Slater interface; global priority for that consequence is not established.

A concrete recent comparator is Kang, Kwon, Scarola, and Park \cite{Kang}: they prepare arbitrary BCS states and implement Gutzwiller projection with amplitude amplification. Their query cost depends on $W^{-1/2}$, where $W$ is projected weight. Theorem~\ref{thm:main} removes that dependence for the explicitly represented number-conserving Slater subclass, under the fresh FPRAS import. It does not extend their generic paired-BCS conclusion.

This advantage need not imply a compact matrix-product state. Take $k$ disjoint pairs, each connecting the first $k$ sites to the last $k$ sites. Occupy two orbitals per pair, one up and one down, each equally supported across its two sites. Single-occupancy projection has probability $2^{-k}$ and produces $k$ singlets. Across the stated cut its Schmidt coefficients are all $2^{-k/2}$. Any state of Schmidt rank at most $D$, including an MPS with bond $D$ in that order, has squared overlap at most $D/2^k$. Trace error $\epsilon$ requires $D\ge(1-\epsilon^2)2^k$. This elementary family witnesses the difference between efficient coherent access and that fixed-order representation.

Synchronized triple products of determinant weights give a different boundary. Three represented matroids can encode the three coordinates of three-dimensional matching, so a common nonzero basis can decide perfect matching. Ohsaka--Matsuoka \cite{OM} establish strong approximation hardness for triple products of principal-minor weights. That projector synchronizes occupied subsets across species and is not a partition quota. It gives no blanket hardness statement for three flavors with one particle per site, which remain inside Theorem~\ref{thm:main}.

For a paired state with skew matrix $F$, amplitudes are $\Pf(F[S])$ on even subsets. Fixed-size Pfaffian weight polynomials need not be log-concave: a choice with two disjoint pairing edges has degree-two polynomial $x_1x_2+x_3x_4$, whose Hessian has two positive eigenvalues. This proves failure of the determinant route, not computational hardness. There is even a different support mechanism. The family of nonzero principal Pfaffians is the represented delta-matroid $D(F)$. For physical pairs, the family of exactly-one-per-pair subsets is a twist, by one fixed element per pair, of the matching delta-matroid. Their intersection is precisely projected support. Represented linear delta-matroid intersection has randomized polynomial field-operation decision and witness algorithms \cite{KW}; a returned witness is directly verifiable by its Pfaffian. For rational-complex inputs, clear denominators and bound the squared norms of all Gaussian-integral principal Pfaffians. Reduction modulo a larger odd prime, embedding $i$ in a degree-at-most-two finite-field extension, preserves every support; the required prime has polynomial bit length. Thus a zero-support obstruction cannot simply be asserted for this class. This support observation supplies neither a weighted norm FPRAS nor coherent preparation. Sequential charge/parity detection interleaved with fermionic linear optics \cite{Beenakker,DT} is a different computational model and does not by itself prove hardness of a single projected-BCS layer.

The conditional status is material. We manually examined the source's load-bearing counting chain (printed pp.~18--25) and rational bit argument (p.~38), including its overlap penalty, tilted-field complement bound, and conditional restrictions. We found no obvious gap in those inspected arguments. We did not independently verify every theorem in the 65-page manuscript, run its FPRAS, implement this quantum compiler, or obtain external specialist validation. Exact small-instance diagnostics check determinant coefficient identities, complex local-estimator identities, rare-support examples, and conditional-error inequalities; they are not execution evidence for the imported algorithm. Independent internal audits agree with the transfer proof. A full source audit, implementation, and comprehensive priority review remain required before presenting an unconditional or externally validated result.

\begin{thebibliography}{99}\small
\bibitem{CL} X.~Chen and K.~Liu, \emph{An FPRAS for Counting Common Bases of Two Matroids}, arXiv:2610.06724v1 (5 October 2026). \url{https://arxiv.org/abs/2610.06724}.
\bibitem{GR} L.~Grover and T.~Rudolph, \emph{Creating superpositions that correspond to efficiently integrable probability distributions}, quant-ph/0208112 (2002). \url{https://arxiv.org/abs/quant-ph/0208112}.
\bibitem{Hav} V.~Havl\'{i}\v{c}ek, \emph{Amplitude ratios and neural network quantum states}, Quantum 7, 938 (2023); arXiv:2201.09128. \url{https://doi.org/10.22331/q-2023-03-02-938}.
\bibitem{BH} P.~Br\"and\'en and J.~Huh, \emph{Lorentzian polynomials}, Annals of Mathematics 192 (2020), 821--891. \url{https://doi.org/10.4007/annals.2020.192.3.4}.
\bibitem{ALOV} N.~Anari, K.~Liu, S.~Oveis Gharan, and C.~Vinzant, \emph{Log-concave polynomials II: high-dimensional walks and an FPRAS for counting bases of a matroid}, STOC (2019), 1--12. \url{https://doi.org/10.1145/3313276.3316385}.
\bibitem{Edmonds} J.~Edmonds, \emph{Matroid intersection}, Annals of Discrete Mathematics 4 (1979), 39--49. \url{https://doi.org/10.1016/S0167-5060(08)70817-3}.
\bibitem{Kang} B.~Kang, H.~Kwon, V.~W.~Scarola, and K.~Park, \emph{Scalable Quantum Algorithms for Gutzwiller Projection}, arXiv:2606.06919 (2026). \url{https://arxiv.org/abs/2606.06919}.
\bibitem{OM} N.~Ohsaka and T.~Matsuoka, \emph{Computational Complexity of Normalizing Constants for the Product of Determinantal Point Processes}, arXiv:2111.14148. \url{https://arxiv.org/abs/2111.14148}.
\bibitem{KW} T.~Koana and M.~Wahlstr\"om, \emph{Faster algorithms on linear delta-matroids}, STACS (2025), 62:1--62:19. \url{https://doi.org/10.4230/LIPIcs.STACS.2025.62}.
\bibitem{Beenakker} C.~W.~J.~Beenakker, D.~P.~DiVincenzo, C.~Emary, and M.~Kindermann, \emph{Charge detection enables free-electron quantum computation}, quant-ph/0401066 (2004). \url{https://arxiv.org/abs/quant-ph/0401066}.
\bibitem{DT} D.~P.~DiVincenzo and B.~M.~Terhal, \emph{Fermionic linear optics revisited}, quant-ph/0403031 (2004). \url{https://arxiv.org/abs/quant-ph/0403031}.
\end{thebibliography}
\end{document}
